% Section: a normalisation convention creates the untrained advantage. All numbers via macros.
\subsection{Decomposition in THINGS}
\label{sec:decomposition}
Several candidate explanations of the THINGS-720 advantage were tested and excluded before the
convention was identified (Appendix~\ref{app:why}). Sampling: subsets of 100 of the 720 images kept a
contrast above the test-image value in \RuleWOne\% of 1000 subsets (W1 met), so the absence of the
advantage in the test images is not sampling variation of 100 images. Measurement noise: the
single-presentation test-image contrast was not positive either (N3 not met). Presentation order:
partialling the position RDM changed $\Delta$ by \RuleWFive{} (W5 met). Carry-over: partialling the
luminance of the preceding image changed it by $\RuleCoOne$ (CO1 not met). Image statistics explained
only part: after partialling global luminance and contrast $\Delta$ was $\RuleWTwo$ (W2 not met), and
after adding coarse luminance $\RuleWSix$ (W6 not met; \RuleWSixFraction\% of the contrast explained;
interpretation of both steps: ``not identified''). Run centring reduced $\Delta$ to $\RuleRcOne$ (RC1
met), but the further predictions of that step were not met (RC2, RC3, RC4), its interpretation rule
had no case for this pattern (``no statement of the rule applies''), and an alignment of run-level drift
with the untrained network's features was not supported (X1 not met; interpretation: ``not identified'').

The decomposition of run centring (step 47c; Table~\ref{tab:decomp}) located the effect in the
session-mean profile $m$. Removing $m$ alone reduced $\Delta$ from $\DecVoneVZero$ to $\DecVoneVOne$
(V1), whereas removing the run deviations alone left $\DecVoneVTwo$ (V2; Z1 not met), and re-adding
either part of the run deviations to the run-centred patterns left it small ($\DecVoneVFour$ and
$\DecVoneVFive$; Z2 not met). The interpretation rule gave: ``mainly removal of the session mean profile
(equivalent to voxel centring), not between-run structure; the description of step 47 as run structure is
to be corrected in the manuscript.'' This agrees with the multiverse of version~1, in which centring the
voxels over the 720 images (raw / correlation / center720) gave $\CenterSevenTwentyCell$
[$\CenterSevenTwentyCellLo$, $\CenterSevenTwentyCellHi$]: the 720 images lie in one session, and
centring over them removes nearly the same profile.

\subsection{Conventions across data sets}
\label{sec:conventions}
% generated by scripts/tables_v2.py; every number is a macro of numbers.tex
\begin{table}[t]
\centering
\small
\begin{adjustbox}{max width=\textwidth, max totalheight=0.82\textheight}
\begin{tabular}{@{}llrrr@{}}
\toprule
Data & Convention & raw / correlation & z / correlation & raw / Euclidean \\
\midrule
THINGS-720 & T-raw & $\CvTsevTrawVoneRawCorr$ (\CvTsevTrawVoneRawCorrPos/3) & $\CvTsevTrawVoneZCorr$ (\CvTsevTrawVoneZCorrPos/3) & $\CvTsevTrawVoneRawEucl$ (\CvTsevTrawVoneRawEuclPos/3) \\
THINGS-720 & T-global & $\CvTsevTglobalVoneRawCorr$ (\CvTsevTglobalVoneRawCorrPos/3) & $\CvTsevTglobalVoneZCorr$ (\CvTsevTglobalVoneZCorrPos/3) & $\CvTsevTglobalVoneRawEucl$ (\CvTsevTglobalVoneRawEuclPos/3) \\
THINGS-720 & T-session & $\CvTsevTsessionVoneRawCorr$ (\CvTsevTsessionVoneRawCorrPos/3) & $\CvTsevTsessionVoneZCorr$ (\CvTsevTsessionVoneZCorrPos/3) & $\CvTsevTsessionVoneRawEucl$ (\CvTsevTsessionVoneRawEuclPos/3) \\
\addlinespace
THINGS test, $k = 1$ & T-raw & $\CvTtestTrawVoneRawCorrKOne$ (\CvTtestTrawVoneRawCorrKOnePos/3) & $\CvTtestTrawVoneZCorrKOne$ (\CvTtestTrawVoneZCorrKOnePos/3) & $\CvTtestTrawVoneRawEuclKOne$ (\CvTtestTrawVoneRawEuclKOnePos/3) \\
THINGS test, $k = 1$ & T-global & $\CvTtestTglobalVoneRawCorrKOne$ (\CvTtestTglobalVoneRawCorrKOnePos/3) & $\CvTtestTglobalVoneZCorrKOne$ (\CvTtestTglobalVoneZCorrKOnePos/3) & $\CvTtestTglobalVoneRawEuclKOne$ (\CvTtestTglobalVoneRawEuclKOnePos/3) \\
THINGS test, $k = 1$ & T-session & $\CvTtestTsessionVoneRawCorrKOne$ (\CvTtestTsessionVoneRawCorrKOnePos/3) & $\CvTtestTsessionVoneZCorrKOne$ (\CvTtestTsessionVoneZCorrKOnePos/3) & $\CvTtestTsessionVoneRawEuclKOne$ (\CvTtestTsessionVoneRawEuclKOnePos/3) \\
\addlinespace
THINGS test, $k = 12$ & T-raw & $\CvTtestTrawVoneRawCorrKOneTwo$ (\CvTtestTrawVoneRawCorrKOneTwoPos/3) & $\CvTtestTrawVoneZCorrKOneTwo$ (\CvTtestTrawVoneZCorrKOneTwoPos/3) & $\CvTtestTrawVoneRawEuclKOneTwo$ (\CvTtestTrawVoneRawEuclKOneTwoPos/3) \\
THINGS test, $k = 12$ & T-global & $\CvTtestTglobalVoneRawCorrKOneTwo$ (\CvTtestTglobalVoneRawCorrKOneTwoPos/3) & $\CvTtestTglobalVoneZCorrKOneTwo$ (\CvTtestTglobalVoneZCorrKOneTwoPos/3) & $\CvTtestTglobalVoneRawEuclKOneTwo$ (\CvTtestTglobalVoneRawEuclKOneTwoPos/3) \\
THINGS test, $k = 12$ & T-session & $\CvTtestTsessionVoneRawCorrKOneTwo$ (\CvTtestTsessionVoneRawCorrKOneTwoPos/3) & $\CvTtestTsessionVoneZCorrKOneTwo$ (\CvTtestTsessionVoneZCorrKOneTwoPos/3) & $\CvTtestTsessionVoneRawEuclKOneTwo$ (\CvTtestTsessionVoneRawEuclKOneTwoPos/3) \\
\addlinespace
NSD, $k = 1$ & N-raw & $\CvNsdNrawVoneRawCorrKOne$ (\CvNsdNrawVoneRawCorrKOnePos/8) & $\CvNsdNrawVoneZCorrKOne$ (\CvNsdNrawVoneZCorrKOnePos/8) & $\CvNsdNrawVoneRawEuclKOne$ (\CvNsdNrawVoneRawEuclKOnePos/8) \\
NSD, $k = 1$ & N-global & $\CvNsdNglobalVoneRawCorrKOne$ (\CvNsdNglobalVoneRawCorrKOnePos/8) & $\CvNsdNglobalVoneZCorrKOne$ (\CvNsdNglobalVoneZCorrKOnePos/8) & $\CvNsdNglobalVoneRawEuclKOne$ (\CvNsdNglobalVoneRawEuclKOnePos/8) \\
NSD, $k = 1$ & N-session & $\CvNsdNsessionVoneRawCorrKOne$ (\CvNsdNsessionVoneRawCorrKOnePos/8) & $\CvNsdNsessionVoneZCorrKOne$ (\CvNsdNsessionVoneZCorrKOnePos/8) & $\CvNsdNsessionVoneRawEuclKOne$ (\CvNsdNsessionVoneRawEuclKOnePos/8) \\
\addlinespace
NSD, $k = 3$ & N-raw & $\CvNsdNrawVoneRawCorrKThree$ (\CvNsdNrawVoneRawCorrKThreePos/8) & $\CvNsdNrawVoneZCorrKThree$ (\CvNsdNrawVoneZCorrKThreePos/8) & $\CvNsdNrawVoneRawEuclKThree$ (\CvNsdNrawVoneRawEuclKThreePos/8) \\
NSD, $k = 3$ & N-global & $\CvNsdNglobalVoneRawCorrKThree$ (\CvNsdNglobalVoneRawCorrKThreePos/8) & $\CvNsdNglobalVoneZCorrKThree$ (\CvNsdNglobalVoneZCorrKThreePos/8) & $\CvNsdNglobalVoneRawEuclKThree$ (\CvNsdNglobalVoneRawEuclKThreePos/8) \\
NSD, $k = 3$ & N-session & $\CvNsdNsessionVoneRawCorrKThree$ (\CvNsdNsessionVoneRawCorrKThreePos/8) & $\CvNsdNsessionVoneZCorrKThree$ (\CvNsdNsessionVoneZCorrKThreePos/8) & $\CvNsdNsessionVoneRawEuclKThree$ (\CvNsdNsessionVoneRawEuclKThreePos/8) \\
\bottomrule
\end{tabular}
\end{adjustbox}
\caption{Untrained minus Backprop ($\Delta$) at V1 under the three voxel normalisation conventions
(Section~\ref{sec:convmethods}), for the paper cell (raw model units, correlation distance), the
standardised cell (model units z-scored over the images) and raw units with Euclidean distance.
Mean over seeds 0--4, then subjects; in brackets the number of subjects with $\Delta > 0$.
T-global and N-session are the conventions of the respective data sets (THINGS: the paper; NSD:
the NSD standard).}
\label{tab:conventions}
\end{table}



Table~\ref{tab:conventions} and Fig.~\ref{fig:conv} show the paper cell under the three conventions of
each data set (step 48). In THINGS-720 the advantage grew when the betas were not normalised
($\CvTsevTrawVoneRawCorr$ under T-raw) and nearly vanished when they were normalised within session
($\CvTsevTsessionVoneRawCorr$ under T-session; NB3 met). The THINGS test images, whose presentations are
spread over 12 sessions, showed no advantage under T-global ($\CvTtestTglobalVoneRawCorrKOne$ at $k = 1$)
but a clear one under T-raw ($\CvTtestTrawVoneRawCorrKOne$ at $k = 1$, $\CvTtestTrawVoneRawCorrKOneTwo$ at
$k = 12$). Normalising over all trials leaves a profile common to the analysed trials only when they come
from one session: under N-global, which normalises NSD over all sessions as T-global does for THINGS, the
paper cell was $\CvNsdNglobalVoneRawCorrKOne$ at $k = 1$ ($\CvNsdNglobalVoneRawCorrKThree$ at $k = 3$;
\CvNsdNglobalVoneRawCorrKOnePos{} of 8 subjects above 0 in both),
because the presentations of the 515 images are spread over 30 sessions; the THINGS test images, spread
over 12 sessions, likewise showed no advantage under T-global. In NSD, the pre-specified prediction that raw betas create the advantage was met: under N-raw
the paper cell was $\CvNsdNrawVoneRawCorrKOne$ at $k = 1$, positive in \CvNsdNrawVoneRawCorrKOnePos{} of
8 subjects (sign-flip $p = \CvNsdNrawVoneRawCorrKOneP$), against $\CvNsdNsessionVoneRawCorrKOne$ under the
NSD standard (NB1 and NB2 met). Under N-raw the multiverse had \MvNrawAbove{} distinct analysis with an
interval above 0, the paper cell itself; \MvNrawBelow{} were below 0 and \MvNrawIncl{} included 0
(Fig.~\ref{fig:multi}). Standardising the model units removed most of the effect
($\CvTsevTrawVoneZCorr$ under T-raw, $\CvNsdNrawVoneZCorrKOne$ under N-raw), and Euclidean distance was
almost unaffected by the convention. The untrained advantage of version~1 is therefore a consequence of
the normalisation convention: it appears when a response profile common to all trials is left in the
voxel patterns and the model units are not centred, under correlation distance.

\subsection{Response amplitude}
\label{sec:amplitude}
With a common profile left in, the correlation distance between two patterns depends on how strongly
each is aligned with the profile. In THINGS-720 under T-global, this brain amplitude $g$ correlated more
strongly with contrast energy than with luminance in every subject (G1 met; Spearman with CE
$\AmpTsevTglobalVoneSubOneCE$, $\AmpTsevTglobalVoneSubTwoCE$, $\AmpTsevTglobalVoneSubThreeCE$; with $L$
$\AmpTsevTglobalVoneSubOneL$, $\AmpTsevTglobalVoneSubTwoL$, $\AmpTsevTglobalVoneSubThreeL$), consistent
with the dependence of V1 responses (Fig.~\ref{fig:amp}a) to natural images on contrast energy \citep{kay2013}. Both networks
matched this amplitude about equally: the Spearman correlation between $g$ and the model amplitude $g^M$
was larger for the untrained network in only \RuleGTwo{} of 3 THINGS subjects (G2 not met; mean
$\AmpGmTsevTglobalVoneKOneRand$ against $\AmpGmTsevTglobalVoneKOneBp$) and \RuleGThree{} of 8 NSD
subjects under N-raw (G3 not met; $\AmpGmNsdNrawVoneKOneRand$ against $\AmpGmNsdNrawVoneKOneBp$). The
interpretation rule of step 48 has no case for this pattern (NB1 and NB2 met, G2 or G3 not met), and its
outcome is: ``no statement of the rule applies''.

Step 48b asked instead how much each RDM is dominated by amplitude structure, using the sum form
$A_{ij} = 2 - g_i - g_j$ (Section~\ref{sec:ampmethods}). The untrained network's RDM was more
amplitude-dominated than Backprop's in every seed in each data set (D1 met; Fig.~\ref{fig:amp}b; mean Spearman(RDM, $A^M$)
$\DomTsevConvOneSumRand$ against $\DomTsevConvOneSumBp$ in THINGS-720, $\DomTestConvOneSumRand$ against
$\DomTestConvOneSumBp$ in the test images, $\DomNsdConvOneSumRand$ against $\DomNsdConvOneSumBp$ in NSD).
The brain RDM was amplitude-dominated when the profile was left in and not when it was removed (D2 met;
THINGS-720 $\BDomTsevTrawVoneKOneSum$ under T-raw against $\BDomTsevTsessionVoneKOneSum$ under T-session;
NSD $\BDomNsdNrawVoneKOneSum$ under N-raw against $\BDomNsdNsessionVoneKOneSum$ under N-session, larger in
\RuleDTwoNsd{} of 8 subjects). Removing the amplitude structure from both RDMs by the two-sided partial
reduced the NSD advantage under N-raw from $\PbNsdNrawVoneKOneNone$ to $\PbNsdNrawVoneKOneSum$, with
$\Delta \leq +0.020$ in \RuleDThreeLe{} of 8 subjects (D3 met; Table~\ref{tab:partial48b}). The
interpretation rule gave: ``the untrained advantage under conventions that leave a common profile in both
patterns is attributed to amplitude structure dominating both RDMs, with the untrained network's RDM more
amplitude-dominated; not to a better match of per-image amplitude.'' Descriptively, the model's amplitude
dominance followed the moment ratio $m$ of version~1 across the 10 checkpoints (Spearman
$\DomMrNsdConvOneSum$ in NSD and $\DomMrTsevConvOneSum$ in THINGS-720; $m$ $\MratioTsevConvOneRand$ for the
untrained network against $\MratioTsevConvOneBp$ for Backprop in THINGS-720), which links the result to the
mean-profile alignment of Section~\ref{sec:gainaxis}. With the difference form $|g_i - g_j|$ the partial
reduced the NSD advantage only to $\PbNsdNrawVoneKOneDiff$ (Appendix~\ref{app:conv}).

\subsection{Why ridge is unaffected}
\label{sec:ridgeunaffected}
A response profile common to all trials is a constant per voxel. A voxel-wise readout fits an intercept
per voxel and is scored by a correlation over images per voxel, so the constant is absorbed and the score
does not change. This is consistent with the ridge results under every normalisation that was tested
(RD1, RD2 in NSD; Y1 to Y4 with run centring and without voxel selection). The readout and the
correlation-distance RDM thus differ in exactly the quantity that the convention controls.

\subsection{Open points}
\label{sec:openres}
Three observations are not explained. (1)~In THINGS-720 the untrained network was ahead under raw units
with Euclidean distance under every convention ($\CvTsevTrawVoneRawEucl$, $\CvTsevTglobalVoneRawEucl$,
$\CvTsevTsessionVoneRawEucl$, each in \CvTsevTglobalVoneRawEuclPos{} of 3 subjects). Euclidean distance is
unaffected by subtracting a profile common to all patterns, so this contrast is not created by the
convention; it is absent in NSD (Table~\ref{tab:convfull}). (2)~After removing the amplitude structure
from both RDMs, a residual untrained advantage remains in THINGS, unlike NSD: $\PbTsevTrawVoneKOneSum$
under T-raw (\PbTsevTrawVoneKOneSumPos{} of 3 subjects positive), $\PbTsevTglobalVoneKOneSum$ under
T-global, and $\PbTestTrawVoneKOneTwoSum$ in the test images under T-raw at $k = 12$. (3)~The sign of the
relation between the brain amplitude and the image statistics depends on the convention: in NSD the
Spearman correlation of $g$ with contrast energy was $\AmpNsdNrawVoneKOneCE$ under N-raw and
$\AmpNsdNsessionVoneKOneCE$ under N-session, and in the test images $\AmpTestTrawVoneKOneTwoCE$ under T-raw
and $\AmpTestTglobalVoneKOneTwoCE$ under T-global (Appendix~\ref{app:conv}, Table~\ref{tab:amplitude}).

\subsection{A side finding: presentation order}
\label{sec:sidefinding}
The V1 RDM of sub-02 correlated with the presentation-position RDM (Spearman $\OrderRhoSubTwo$; sub-01
$\OrderRhoSubOne$, sub-03 $\OrderRhoSubThree$). sub-01 and sub-02 saw the 720 images in the same order and
run assignment, so their RDMs share any order-related component, and the agreement between these two
subjects in THINGS is a limited test of cross-subject consistency. Partialling the position RDM did not
change the group contrast (W5).
